Bayesian inference requires evaluating the likelihood function at
thousands to millions of points in parameter space. When each
likelihood evaluation demands a forward-model simulation that takes
seconds to minutes, direct sampling methods such as Markov chain Monte
Carlo (MCMC) become prohibitively expensive. To overcome this
computational bottleneck, \texttt{alabi} constructs a \emph{surrogate
model}: a fast, continuously differentiable approximation to the
log-likelihood surface that can be evaluated in microseconds once
trained \citep{Luger2017,Fleming2024}.

The surrogate model is a Gaussian process (GP), a non-parametric
regression method that defines a probability distribution over
functions rather than fitting a fixed functional form
\citep{Rasmussen2006}. Given a set of training points---pairs of
input parameters and their corresponding log-likelihood values---the
GP predicts the log-likelihood at any untested point together with an
uncertainty estimate. This uncertainty quantification is a key
advantage over parametric interpolation schemes: it reveals where the
surrogate is poorly constrained and thereby guides the acquisition of
new training points during the active learning phase described in
Section~\ref{sec:active_learning}.

A GP is fully specified by two components: a mean function and a
covariance function, the latter commonly called a \emph{kernel}. The
mean function captures the global trend of the target function (often
set to a constant), while the kernel controls how the GP generalizes
from observed points to unobserved ones. Mathematically, the kernel is a function
$k(\mathbf{x}, \mathbf{x}')$ that returns the covariance between the
GP's predictions at any two points $\mathbf{x}$ and $\mathbf{x}'$ in
parameter space. When two points are close together (relative to the
length scale $\ell$), the kernel returns a value near the output
variance $\sigma^{2}$, meaning the GP expects their log-likelihood
values to be strongly correlated. As the points move farther apart,
the kernel value decays toward zero and the predictions become
independent. The functional form of the kernel
therefore encodes prior assumptions about the smoothness and
correlation structure of the log-likelihood surface. Choosing a kernel
that is too smooth can wash out genuine features, while a kernel that
is too rough can overfit noise in the training data. Each kernel has a small number of free parameters---an output
variance $\sigma^{2}$ and one or more length scales $\ell$---that
must be tuned to the training data. The standard approach is to
maximize the \emph{marginal likelihood} (also called the model
evidence), which is the probability of the observed training outputs
given the inputs after integrating over all possible functions
consistent with the GP prior. The marginal likelihood naturally
balances model fit against model complexity: a kernel configuration
that fits the training data well but requires an implausibly flexible
function space is penalized relative to a simpler configuration that
fits nearly as well \citep{Rasmussen2006}. In practice,
\texttt{alabi} maximizes the log of the marginal likelihood using the
L-BFGS-B gradient-based optimizer \citep{Byrd1995}, starting from
multiple initial guesses to reduce the chance of converging to a local
maximum.

In addition to the kernel, GP performance depends on how the input
parameters and output log-likelihood values are numerically
represented. The input parameters in a planetary-evolution model often
span very different physical scales (e.g., masses in Earth masses
versus orbital periods in days), and the log-likelihood values can
range over many orders of magnitude. Without preprocessing, parameters
with large numerical ranges can dominate the distance calculations in
the kernel, effectively rendering the GP insensitive to parameters
with small ranges. Similarly, extreme dynamic range in the output can
cause the GP to focus its predictive power on the largest-magnitude
training points at the expense of accuracy near the likelihood peak,
which is precisely the region that matters most for inference.
\emph{Scaling}---applying a mathematical transformation to the inputs,
the outputs, or both before training---mitigates these problems by
placing all quantities on comparable numerical footing.

Because the optimal kernel and scaling strategy depend on the specific
shape of the log-likelihood surface, which is not known \emph{a
priori}, we perform a grid search over three kernel families and
several data-scaling strategies. The goal of this search is to
identify the kernel--scaler combination that most accurately
reproduces the true log-likelihood surface, as measured by the
prediction error on a held-out test set. The winning configuration is
then used for the full surrogate-model training.

We consider three stationary kernels. The squared exponential (SE)
kernel,
\begin{equation}
  k_{\mathrm{SE}}(\mathbf{x}, \mathbf{x}') = \sigma^{2}
  \exp\!\left(-\frac{|\mathbf{x} - \mathbf{x}'|^{2}}{2\ell^{2}}\right),
  \label{eq:se_kernel}
\end{equation}
assumes infinitely differentiable target functions and produces the
smoothest interpolants. The Mat\'{e}rn 3/2 kernel,
\begin{equation}
  k_{3/2}(\mathbf{x}, \mathbf{x}') = \sigma^{2}
  \left(1 + \frac{\sqrt{3}\,|\mathbf{x} - \mathbf{x}'|}{\ell}\right)
  \exp\!\left(-\frac{\sqrt{3}\,|\mathbf{x} - \mathbf{x}'|}{\ell}\right),
  \label{eq:matern32_kernel}
\end{equation}
produces once-differentiable sample paths and accommodates rougher
target functions. The Mat\'{e}rn 5/2 kernel,
\begin{equation}
  k_{5/2}(\mathbf{x}, \mathbf{x}') = \sigma^{2}
  \left(1 + \frac{\sqrt{5}\,|\mathbf{x} - \mathbf{x}'|}{\ell}
  + \frac{5\,|\mathbf{x} - \mathbf{x}'|^{2}}{3\ell^{2}}\right)
  \exp\!\left(-\frac{\sqrt{5}\,|\mathbf{x} - \mathbf{x}'|}{\ell}\right),
  \label{eq:matern52_kernel}
\end{equation}
is twice-differentiable and occupies the middle ground between the SE
and Mat\'{e}rn 3/2 kernels. In each case, $\sigma^{2}$ is the output
variance and $\ell$ is the characteristic length scale.

In addition to the kernel, the performance of a GP depends on how
the input (parameter) and output (log-likelihood) values are scaled
before training. We consider three input scalers: no transformation,
min--max scaling to $[0,1]$, and standardization to zero mean and unit
variance. For the output, we add a fourth option: the negative
logarithm, $y \to -\!\ln(-y)$, which compresses the dynamic range of
strongly negative log-likelihoods. The full grid therefore comprises
$3 \times 3 \times 4 = 36$ kernel--scaler combinations.

For each combination, the GP is trained on $N_{\mathrm{train}} = 500$
points drawn from the five-dimensional parameter space via Latin
hypercube sampling (LHS) and evaluated on a separate set of
$N_{\mathrm{test}} = 500$ LHS points. The quality of each
configuration is quantified by the mean squared error (MSE) between
the GP predictions and the true log-likelihood values on the test set,
\begin{equation}
  \mathrm{MSE} = \frac{1}{N_{\mathrm{test}}}
  \sum_{i=1}^{N_{\mathrm{test}}}
  \bigl[\hat{y}(\mathbf{x}_{i}) - y(\mathbf{x}_{i})\bigr]^{2},
  \label{eq:mse}
\end{equation}
where $\hat{y}(\mathbf{x}_{i})$ is the GP prediction and
$y(\mathbf{x}_{i})$ is the forward-model log-likelihood at test point
$\mathbf{x}_{i}$. Both quantities are evaluated in the original
(unscaled) output space to ensure a fair comparison across scaling
strategies.

Table~\ref{tab:grid_search} reports the test MSE for all 36
configurations. The best-performing configuration is the SE kernel
with no input scaling and min--max output scaling
($\mathrm{MSE} = 3.44$). The SE and Mat\'{e}rn 5/2 kernels with
min--max output scaling consistently achieve the lowest errors
($\mathrm{MSE} \lesssim 3.8$), while the Mat\'{e}rn 3/2 kernel
performs notably worse ($\mathrm{MSE} \gtrsim 7.7$), indicating that
the log-likelihood surface is sufficiently smooth to benefit from the
higher differentiability of the SE and Mat\'{e}rn 5/2 kernels. The
negative-log output scaler performs poorly across all kernels
($\mathrm{MSE} > 25$), likely because the transformation distorts the
relative spacing of log-likelihood values. Input scaling has minimal
effect for the SE kernel (MSE varies by $< 0.1$) and a modest effect
for Mat\'{e}rn 5/2 (MSE varies by $< 0.04$), suggesting that the GP
hyperparameter optimization compensates for differences in input
scaling in these cases.

After selecting the optimal configuration from the grid search, the
GP is refined through 500 iterations of active learning, in which new
training points are sequentially chosen to maximize the GP's
predictive uncertainty. The active learning phase is described in
Section~\ref{sec:active_learning}.
